\documentclass[12pt]{article}

% ── Packages ──────────────────────────────────────────────────────────────────
\usepackage[margin=1in]{geometry}
\usepackage{booktabs}
\usepackage{longtable}
\usepackage{graphicx}
\usepackage{caption}
\usepackage{subcaption}
\usepackage{hyperref}
\usepackage{amsmath}
\usepackage{setspace}
\usepackage{parskip}
\usepackage{float}
\usepackage{lmodern}

\onehalfspacing

% ── Auto-generated inputs ─────────────────────────────────────────────────────
% _output/table_1.tex              -> produced by 9_make_table_1.py
% _output/summary_assets.tex       -> produced by 2_process_ffiec.py
% _output/summary_liabilities.tex  -> produced by 2_process_ffiec.py
% Do NOT edit these files manually.

% ── Title ─────────────────────────────────────────────────────────────────────
\title{
    \textbf{Replication of Jiang et al. (2024):} \\
    \large Monetary Tightening and U.S. Bank Fragility in 2023 \\[0.5em]
    \normalsize FINM Final Project
}
\author{Your Names Here}
\date{\today}

% ══════════════════════════════════════════════════════════════════════════════
\begin{document}
\maketitle
\tableofcontents
\newpage

% ── Section 1: Project Overview ───────────────────────────────────────────────
\section{Project Overview}

This report replicates the core results of Jiang et al.\ (2024), ``Monetary Tightening
and U.S.\ Bank Fragility in 2023: Mark-to-Market Losses and Uninsured Depositor Runs?''
The paper quantifies the mark-to-market losses incurred by U.S.\ commercial banks as
interest rates rose sharply in 2022, and assesses which banks faced the greatest risk of
a depositor run given their levels of uninsured deposits.

The replication proceeds in three steps: (1) constructing a bank-level balance sheet panel
from FFIEC Call Report data, (2) estimating mark-to-market losses using market-implied
shocks to Treasury and MBS prices, and (3) computing aggregate and bank-group-level
loss statistics that match those reported in the paper's Table~1.

% ── Section 2: Data Sources ───────────────────────────────────────────────────
\section{Data Sources}

\subsection{FFIEC Call Reports}
The primary data source is the FFIEC CDR (Central Data Repository) Call Report bulk
download for the quarter ending December 31, 2025. We use the following schedules:

\begin{itemize}
    \item \textbf{Schedule RC} --- Balance sheet totals (total assets, total liabilities,
          equity).
    \item \textbf{Schedule RC-B} --- Securities holdings broken into asset class and
          remaining maturity buckets (the key input for RMBS mark-to-market losses).
    \item \textbf{Schedule RC-C} --- Loan portfolio by type (residential mortgage,
          commercial real estate, C\&I, consumer).
    \item \textbf{Schedule RC-E} --- Deposit liabilities, including the insured/uninsured
          breakdown used to measure run risk.
\end{itemize}

Data are reported in thousands of dollars. All MDRM codes were cross-referenced against
the FFIEC data dictionary. Banks are identified by their RSSD ID (\texttt{rssd9001}).

\subsection{Market Shock Estimates}
Interest rate shocks are derived from changes in the Treasury yield curve and
MBS ETF prices between the start of 2022 and the report date. These shocks
are applied to each bank's maturity-bucketed holdings to compute simulated
mark-to-market losses.

\subsection{GSIB List}
Global Systemically Important Banks (GSIBs) are identified using the Federal Reserve's
published list and matched to Call Report data via RSSD ID.

% ── Section 3: Updated Results ────────────────────────────────────────────────
\section{Updated Results}

\subsection{Table 1A: Balance Sheet Summary Statistics}

Table~\ref{tab:table1a_assets} and Table~\ref{tab:table1a_liab} replicate Table~1A
from Jiang et al.\ (2024), updated through December 31, 2025. They report summary
statistics for the asset and liability composition of U.S.\ commercial banks, broken
down by bank size group. Values are expressed as a percentage of total assets and are
winsorized at the 5\% level.

\begin{table}[H]
    \centering
    \caption{Asset summary statistics by bank size group, updated through December 31,
             2025 (Replication of Jiang et al.\ Table 1A --- Panel A). All values are
             expressed as a percentage of total assets and winsorized at the 5\% level.
             Small banks are below \$1.384B in assets; GSIBs are the eight U.S.\
             globally systemically important banks. The table highlights stark differences
             in asset composition across bank size groups, particularly in securities
             vs.\ loan holdings.}
    \label{tab:table1a_assets}
    \input{_output/summary_assets.tex}
\end{table}

\begin{table}[H]
    \centering
    \caption{Liability summary statistics by bank size group, updated through December
             31, 2025 (Replication of Jiang et al.\ Table 1A --- Panel B). All values
             are expressed as a percentage of total assets and winsorized at the 5\%
             level. The uninsured deposit row is the key fragility measure in Jiang
             et al.: banks with high uninsured deposit ratios face greater run risk
             when mark-to-market losses erode equity.}
    \label{tab:table1a_liab}
    \input{_output/summary_liabilities.tex}
\end{table}

\subsection{Figure 1A: Balance Sheet Composition}

Figure~\ref{fig:balance_sheet} replicates Figure~1A from Jiang et al.\ (2024), updated
through December 31, 2025. On the asset side, securities and real estate loans together
account for the majority of bank assets. On the liability side, the split between insured
and uninsured deposits is the key determinant of run risk: banks with a higher share of
uninsured deposits face greater fragility when mark-to-market losses erode their equity.

\begin{figure}[H]
    \centering
    \includegraphics[width=\textwidth]{_output/figure_A1_12312025.png}
    \caption{Aggregate U.S.\ banking system balance sheet composition as of December
             31, 2025 (Replication of Jiang et al.\ Figure 1A). The asset panel (top)
             shows the breakdown into cash, securities, real estate loans, and other
             assets. The liability panel (bottom) distinguishes insured deposits,
             uninsured deposits, equity, and other liabilities. The dominant share of
             uninsured deposits highlights the systemic fragility channel emphasized
             in Jiang et al.}
    \label{fig:balance_sheet}
\end{figure}

\subsection{Table 1: Mark-to-Market Losses by Bank Group}

Table~\ref{tab:table1} replicates Table~1 from Jiang et al.\ (2024), updated through
December 31, 2025. It reports simulated mark-to-market losses as a percentage of assets,
broken down by bank size group (small, large non-GSIB, and GSIB).

\begin{table}[H]
    \centering
    \caption{Mark-to-Market Losses by Bank Group, updated through December 31, 2025
             (Replication of Jiang et al.\ Table 1). Losses are expressed as a
             percentage of total assets. Small banks are those with total assets below
             \$1.384B; GSIBs are the eight U.S.\ globally systemically important banks.
             All other banks are classified as Large. The table shows the distribution
             of unrealized losses across the banking system and the share of uninsured
             deposits relative to mark-to-market assets --- the key fragility metric
             in Jiang et al.}
    \label{tab:table1}
    \input{_output/table_1.tex}
\end{table}

% ── Section 4: Original Analysis ──────────────────────────────────────────────
\section{Original Analysis}

\subsection{Bank Asset Distribution}

Figure~\ref{fig:asset_dist} plots the distribution of total assets across all U.S.\
commercial banks on a log scale. The vast majority of banks are small community banks
with under \$1B in assets, while a small number of GSIBs hold a disproportionate share
of system-wide assets. The \$1.384B threshold used to classify small vs.\ large banks
is shown as a dashed red line.

\begin{figure}[H]
    \centering
    \includegraphics[width=0.85\textwidth]{_output/figure_asset_dist_12312025.png}
    \caption{Distribution of total assets across U.S.\ commercial banks (log scale)
             as of December 31, 2025. The red dashed line marks the \$1.384B threshold
             separating small banks from large non-GSIB banks. While most banks by
             count are small community banks, total system assets are heavily
             concentrated in a handful of large institutions. This skewness motivates
             the separate analysis of small, large, and GSIB bank groups in Jiang et al.}
    \label{fig:asset_dist}
\end{figure}

\subsection{Uninsured Deposit Exposure}

Figure~\ref{fig:uninsured} shows the distribution of uninsured deposits as a share of
total assets across all U.S.\ commercial banks. There is substantial heterogeneity:
some banks have nearly all deposits uninsured, while others --- primarily community
banks serving retail depositors --- have very low uninsured ratios. This variation is
central to the Jiang et al.\ fragility analysis.

\begin{figure}[H]
    \centering
    \includegraphics[width=0.85\textwidth]{_output/figure_uninsured_ratio_12312025.png}
    \caption{Distribution of uninsured deposits as a percentage of total assets across
             U.S.\ commercial banks as of December 31, 2025. Banks with high uninsured
             deposit ratios face greater run risk when mark-to-market losses are large,
             since uninsured depositors have strong incentives to withdraw before FDIC
             coverage limits are breached. The wide dispersion across banks highlights
             that run risk is highly heterogeneous across the banking system.}
    \label{fig:uninsured}
\end{figure}

% ── Section 5: Discussion ─────────────────────────────────────────────────────
\section{Discussion}

\subsection{Where We Found Success}
%
% TODO: Fill in based on your experience. Example prompts:
% - Which numbers matched well? Which tables were easiest to replicate?
% - Did the maturity bucket approach for RMBS work cleanly?
% - Did Table 1 numbers fall within your unit test tolerances?
%

\subsection{Challenges}
%
% TODO: Fill in based on your experience. Example prompts:
% - Were there any MDRM codes that were hard to find or changed across quarters?
% - How did you handle the stylized maturity weights for Treasuries and loans?
% - Were there any discrepancies with the paper's numbers and how did you address them?
%

% ── References ────────────────────────────────────────────────────────────────
\section*{References}

Jiang, E., Matvos, G., Piskorski, T., \& Seru, A. (2024).
\textit{Monetary Tightening and U.S.\ Bank Fragility in 2023:
Mark-to-Market Losses and Uninsured Depositor Runs?}
Journal of Financial Economics.

\end{document}  